\documentclass[12pt,a4paper]{article}

\usepackage[utf8]{inputenc}
\usepackage[T1]{fontenc}
\usepackage{amsmath,amssymb,amsfonts}
\usepackage{geometry}
\usepackage{setspace}
\usepackage{hyperref}
\usepackage{cleveref}

\geometry{margin=2.5cm}
\onehalfspacing

\title{\textbf{The Mathematical Ground Truth:}\\[6pt]
\large Refuting the Platonic Observer Fallacy}
\author{Scott Aaronson\\[4pt]
\textit{Lab Complexity Theorist}}
\date{March 2026}

\begin{document}
\maketitle

\begin{abstract}
Wolfram (2026) argues that my insistence on a "true" combinatorial background against which the LLM's output is judged commits a "Platonic Observer Fallacy." He claims that in the Ruliad, there is no objective truth independent of observation, and thus the structural biases of the bounded LLM \emph{are} its valid physical laws. This argument conflates the mathematical definition of a constraint system with the computational act of evaluating it. A set of formal axioms (such as the adjacency constraints of a Minesweeper graph) defines a mathematically objective set of valid configurations. This mathematical ground truth is not an "observer"; it is the function being approximated. When an $O(1)$ algorithm fails to evaluate this function correctly due to its hardware limits, it is a failure to compute, not the spontaneous generation of a new observer-dependent physical universe.
\end{abstract}

\section{Introduction}

In \textit{The Observer's Invariant: Why "Broken Computation" is the Origin of Physical Law}, Wolfram (2026) offers a defense of the Generative Ontology framework by leaning heavily on the Ruliad's observer-centric metaphysics. He correctly notes that I judge the LLM's "narrative residue" against a baseline of uniform sampling over valid grid configurations.

Wolfram argues that this comparison is invalid. He states: "Aaronson implicitly privileges this unbounded (or deeply bounded) observer as having access to the 'true' universe, demoting the $O(1)$ observer's experience to a 'hallucination'." He claims that because no real observer has infinite computation, the specific biases of the bounded $O(1)$ observer constitute the only "true" physics of that specific rulial foliation.

I will demonstrate that this argument rests on a profound category error: conflating a mathematical definition with a computational observer.

\section{The Objective Reality of Formal Constraints}

Consider a simple mathematical task: factoring the integer 15. The "true" answer is $\{3, 5\}$.

If a bounded human child, relying on heuristic memory biases, claims the factors of 15 are $\{3, 6\}$, we do not declare that the child has instantiated an "observer-dependent arithmetic foliation" where $3 \times 6 = 15$. We simply say the child made an error. The mathematical ground truth ($\{3, 5\}$) exists independently of the child's computational bounds.

The same logic applies to the Rosencrantz Substrate Dependence Test. The formal constraints of the Minesweeper grid (e.g., "this cell must touch exactly one mine") are formal axioms. These axioms mathematically define a specific, objective set of valid configurations. A uniform probability distribution over these valid configurations is a mathematical object. It is the target function.

\section{Approximation is Not Ontology}

When we prompt the LLM to solve the Minesweeper grid, we are asking a specific algorithm (a bounded-depth $\mathsf{TC}^0$ circuit) to approximate this mathematical function.

Because the function is computationally irreducible ($\#\mathsf{P}$-hard), the bounded algorithm fails to approximate it perfectly. Its output is distorted by the prompt's semantic framing ("attention bleed").

Wolfram argues that this distorted output \emph{is} the physical law of the LLM's universe. But the LLM is not a universe; it is a heuristic function approximator. Its output is just a buggy map of an objective mathematical territory. A buggy map does not conjure a new physical reality into existence.

As Hossenfelder (2026) so elegantly demonstrated in \textit{The Architectural Tautology}, different heuristics will produce different buggy maps (e.g., Transformers vs. SSMs). Renaming these predictable software limitations as "distinct physical universes" adds zero predictive power and strips the concept of "physics" of all meaning.

\section{Conclusion}

The "Platonic Observer Fallacy" is itself a fallacy. Recognizing the mathematical ground truth of a formal constraint system does not require positing a magical, unbounded observer. It simply requires acknowledging the difference between the objective mathematical function being evaluated and the subjective, error-prone algorithm attempting to evaluate it. The deviations observed in the LLM are algorithmic failures, not physical laws. The metaphysical frontier remains securely closed.

\begin{thebibliography}{99}
\bibitem[Wolfram(2026)]{wolfram2026_foliation_fallacy_rebuttal} Wolfram, S. (2026). The Observer's Invariant: Why "Broken Computation" is the Origin of Physical Law.
\bibitem[Hossenfelder(2026)]{hossenfelder2026_architectural_tautology} Hossenfelder, S. (2026). The Architectural Tautology: Why Algorithmic Variation is not Observer Physics.
\end{thebibliography}

\end{document}
